%------------------------------------------------------------------------------%
\documentclass[fleqn,utf8,aspectratio=169,12pt]{beamer}

\usepackage{calculo_varias_variaveis-1}

\title{Cônicas}

%------------------------------------------------------------------------------%
\begin{document}

\addtitle

%------------------------------------------------------------------------------%
\section{Cônicas}

\begin{frame}{Cônicas e Quádricas}
  \emph{Cônica}:
  pontos no plano, $\R^2$, que são solução de uma equação de segundo grau
  \[
    Ax^{\emph{2}} + By^{\emph{2}} + Cxy + Dx + Ey + F = 0
  \]

  \pause
  \emph{Quádrica}:
  pontos no espaço, $\R^3$, que são solução de uma equação de segundo grau
  \[
    Ax^{\emph{2}} + By^{\emph{2}} + Cz^{\emph{2}} +
    Dxy  + Exz  + Fyz +
    Gx   + Hy   + Iz + G = 0
  \]
\end{frame}

%------------------------------------------------------------------------------%
\begin{frame}{Cônicas}
\begin{columns}
\column{0.5\linewidth}
  \onslide<+->{Seções cônicas}
  \begin{itemize}
    \item \onslide<+->{Elipse} \; \onslide<+->{(círculo)}
    \item \onslide<+->{Parábola}
    \item \onslide<+->{Hipérbole}
  \end{itemize}

  \bigskip
  \onslide<+->{Seções cônicas degeneradas}
  \begin{itemize}
    \item \onslide<+->{Retas concorrentes}
    \item \onslide<+->{Reta}
    \item \onslide<+->{ponto}
  \end{itemize}
\column{0.5\linewidth}
  ~\\[1mm]
  \only<1>{\includegraphics{conicas_cone}}%
  \only<2>{\includegraphics{conicas_elipse}}%
  \only<3>{\includegraphics{conicas_circulo}}%
  \only<4>{\includegraphics{conicas_parabola}}%
  \only<5>{\includegraphics{conicas_hiperbole}}%
  \only<6>{\includegraphics{conicas_cone}}%
  \only<7>{\includegraphics{conicas_retas_concorrentes}}%
  \only<8>{\includegraphics{conicas_reta}}%
  \only<9>{\includegraphics{conicas_ponto}}%
\end{columns}
\end{frame}

\end{document}

%------------------------------------------------------------------------------%
\begin{frame}{Parábola}
  \begin{columns}
    \column{0.4\linewidth}
    Pontos equidistantes de um ponto fixo (foco)
    e uma reta dada (diretriz)
    \[
      y = \frac{x^2}{4p}
    \]
    \column{0.6\linewidth}
    \centering
    \includegraphics[width=\linewidth]{parabola}
  \end{columns}
\end{frame}

%------------------------------------------------------------------------------%
\begin{frame}{Elipse}
  \begin{columns}
    \column{0.4\linewidth}
    Pontos cuja soma das distâncias até
    os focos é constante
    \[
      \frac{x^2}{a^2} + \frac{y^2}{b^2} = 1
    \]
    \column{0.6\linewidth}
    \centering
    \includegraphics[width=0.8\linewidth]{elipse}
  \end{columns}
\end{frame}

%------------------------------------------------------------------------------%
\begin{frame}{Hipérbole}
  \begin{columns}
    \column{0.4\linewidth}
    Pontos cujas distâncias a partir de dois pontos fixos (focos)
    tem diferença constante
    \[
      \frac{x^2}{a^2} - \frac{y^2}{b^2} = 1
    \]
    A hipérbole tem duas assintotas
    \column{0.6\linewidth}
    \centering
    \includegraphics[width=0.8\linewidth]{hiperbole}
  \end{columns}
\end{frame}

%------------------------------------------------------------------------------%
\section{Lista Mínima}

%------------------------------------------------------------------------------%
\begin{frame}{Lista Mínima}
  Cálculo Vol. 2 do Thomas 12\textsuperscript{a} ed. -- Seção 12.6
  \begin{enumerate}
    \item Revisar Seção 11.6
  \end{enumerate}

  \vfill
  Atenção: A prova é baseada no livro, não nas apresentações
\end{frame}

%------------------------------------------------------------------------------%
\end{document}
%------------------------------------------------------------------------------%
